\section{A complete analytic endpoint}
\label{rh-analytic}

\subsection{The detecting transform}

\begin{lemma}[Mellin transform of the corrected sum]
\label{rh-mellin}
For $\Re s>1$,
\begin{equation}\label{rh-mellin-formula}
 \int_1^\infty G(x)x^{-s-1}\,dx
 =-\frac{\zeta'(s)/\zeta(s)}{s(s-1)}
       -\frac1{(s-1)^2}+\frac{1+\gamma}{s-1}.
\end{equation}
The right side is holomorphic at $s=1$. At a zero $\rho$ of $\zeta$
in the open critical strip, of multiplicity $m_\rho$, it has residue
\begin{equation}\label{rh-zero-residue}
                       -\frac{m_\rho}{\rho(\rho-1)}\ne0.
\end{equation}
\end{lemma}

\begin{proof}
For $\Re s>1$,
\[
 \int_n^\infty(x/n-1)x^{-s-1}\,dx
       =\frac{n^{-s}}{s(s-1)}.
\]
The series of absolute integrals converges, since
$\sum_n\Lambda(n)n^{-\Re s}<\infty$. Termwise integration and the
logarithmic derivative of the Euler product give the first term in
\eqref{rh-mellin-formula}; integrating the two correcting functions
gives the remaining terms. Write $z=s-1$. The expansion
$\zeta'/\zeta=-z^{-1}+\gamma+O(z)$ shows that the first term is
$z^{-2}-(1+\gamma)z^{-1}+O(1)$, proving removability at one.
At a zero, $\zeta'/\zeta=m_\rho/(s-\rho)+O(1)$, which proves
\eqref{rh-zero-residue}.
\end{proof}

The same calculation for a nonprincipal primitive character gives
\begin{equation}\label{rh-dirichlet-mellin}
 \int_1^\infty G_\chi(x)x^{-s-1}\,dx
 =-\frac{L'(s,\chi)/L(s,\chi)}{s(s-1)}
                 +\frac{L'(1,\chi)/L(1,\chi)}{s-1}.
\end{equation}
The pole at one again cancels. At a zero in the open strip its residue
is \eqref{rh-zero-residue}, with the corresponding multiplicity.

\subsection{Positivity at the real convergence boundary}

\begin{lemma}[A positive Mellin integral]
\label{rh-landau}
Let $f\ge0$ be locally integrable on $[X,\infty)$, $X\ge1$, and
suppose its Mellin integral converges at all sufficiently large real
arguments. If that integral has a holomorphic continuation to a
neighborhood of every real $s>\theta$, then its abscissa of
convergence is at most $\theta$.
\end{lemma}

\begin{proof}
The assertion is immediate if the abscissa is $-\infty$. Otherwise
suppose it is a finite number $\sigma_c>\theta$. Holomorphy near
$\sigma_c$ supplies a real $a>\sigma_c$ at which the continued
function has a Taylor disc of radius $R>a-\sigma_c$.
Differentiation of the integral at $a$ is justified by convergence at
some smaller real argument still greater than $\sigma_c$; each fixed
power of $\log x$ is absorbed by the resulting power of $x$.
For $0<t<R$ the Taylor series is
\[
 F(a-t)=\sum_{k=0}^\infty\frac{t^k}{k!}
              \int_X^\infty f(x)(\log x)^kx^{-a-1}\,dx.
\]
All terms are nonnegative. Tonelli's theorem identifies this convergent
series with
\[
                         \int_X^\infty f(x)x^{-(a-t)-1}\,dx.
\]
Taking $a-\sigma_c<t<R$ contradicts the definition of $\sigma_c$.
\end{proof}

\begin{theorem}[An eventual one-sided criterion for RH]
\label{rh-eventual}
If \eqref{rh-simple-target} holds for every $\epsilon>0$, the
Riemann hypothesis follows.
\end{theorem}

\begin{proof}
Fix $\theta=1/2+\epsilon<1$ and set
$f(x)=C_\epsilon x^\theta-G(x)$ for $x\ge X_\epsilon$.
It is nonnegative. The elementary bound
$|G(x)|\ll x(1+\log^2x)$ gives initial convergence of its Mellin
integral. Formula~\eqref{rh-mellin-formula}, with an entire finite
interval integral subtracted, continues that transform meromorphically.

There is no real singularity above $\theta$. For $s>1$, the Euler
product gives $\zeta(s)>0$. For $0<s<1$, the convergent alternating
series
\[
 \eta(s)=\sum_{n\ge1}(-1)^{n-1}n^{-s}
        =\sum_{k\ge1}\bigl((2k-1)^{-s}-(2k)^{-s}\bigr)>0
\]
satisfies $\zeta(s)=\eta(s)/(1-2^{1-s})<0$. At one the apparent
singularity was removed in Lemma~\ref{rh-mellin}. The extra transform
$C_\epsilon X_\epsilon^{\theta-s}/(s-\theta)$ has its only pole at
$\theta$. Lemma~\ref{rh-landau} therefore gives convergence on
$\Re s>\theta$.

This nonnegative integral is holomorphic there: convergence at a
slightly smaller real argument supplies locally uniform domination,
also after differentiation. Its meromorphic continuation agrees with
that holomorphic function by the identity theorem. A zero of $\zeta$
in the strip with real part greater than $\theta$ would, however,
produce the nonzero residue \eqref{rh-zero-residue}. This is impossible.
Let $\epsilon$ tend to zero. The functional equation reflects a zero
to the left of one half into a zero to the right, completing the proof.
\end{proof}

More generally, it is enough to have eventual estimates
$G(x)\le M_jx^{1/2+\delta_j}$ with $\delta_j\downarrow0$ and finite
$M_j$ depending on $j$. The multiplicative constants need not tend
to one. For any desired $\epsilon>0$, choose $j$ with
$\delta_j<\epsilon$ and increase the initial threshold if necessary.

\subsection{A stronger criterion on arbitrarily large intervals}

We use the following more general analytic theorem from the GRH
manuscript \cite[section ``One-sided criteria for zeros'', proposition
on positive majorants on long intervals]{rh-companion}. Its full
proof is retained there; it is stated here as an imported theorem,
separately from the short proof above.

\begin{theorem}[A meromorphic transform and long intervals]
\label{rh-window-input}
Let $g:[0,\infty)\to\mathbb R$ be measurable with
$|g(u)|\le C_0e^{2u}$ for all $u\ge0$ and some finite $C_0$.
Suppose its Laplace transform has a meromorphic
continuation to $\Re s>0$, with only simple poles, and every real pole
in that half-plane has positive residue. Fix $0<\theta<1$ and $\beta>1$.
Suppose that, for every $\eta>0$, there are arbitrarily large $U$
and measurable nonnegative functions $R_U$ on $[U,\beta U]$ with
\[
 g(u)\le R_U(u)\quad\hbox{almost everywhere},\qquad
 \int_U^{\beta U}R_U(u)e^{-\theta u}\,du\le e^{\eta U}.
\]
Then the transform has no pole in $\Re s>\theta$.
\end{theorem}

\begin{lemma}[Fixed averaging preserves detecting poles]
\label{rh-averaging}
For $g$ satisfying the hypotheses of Theorem~\ref{rh-window-input},
let $g_h(u)=h^{-1}\int_0^hg(u+v)\,dv$ for $h>0$, and let $\Phi$
be the Laplace transform of $g$. Then
\[
 \Phi_h(s)=\frac{e^{hs}-1}{hs}\Phi(s)-E_h(s),\qquad
 E_h(s)=\frac1h\int_0^he^{sv}\int_0^vg(w)e^{-sw}\,dw\,dv.
\]
Here $E_h$ is entire. The multiplier has no zero in $\Re s>0$ and
is positive for real $s>0$.
\end{lemma}

\begin{proof}
Apply Fubini first where the Laplace integral converges absolutely,
and substitute $w=u+v$. Both integrals defining $E_h$ are over compact
intervals, proving entire dependence. The zeros of $e^{hs}-1$ lie
on $\Re s=0$; its value divided by $hs$ at zero is interpreted as one.
\end{proof}

\begin{corollary}[A weighted positive-part criterion for RH]
\label{rh-averaged}
For each $j$, fix $\beta_j>1$, $h_j\ge0$ and
$1/2<\theta_j<1$, with $\theta_j\to1/2$. If
\begin{equation}\label{rh-average-target}
 \liminf_{Y\to\infty}
 \frac{\log\left(1+\displaystyle\int_Y^{Y^{\beta_j}}
       (\mathcal A_{h_j}G(x))_+x^{-\theta_j}\,\frac{dx}{x}\right)}
      {\log Y}=0
\end{equation}
for every $j$, then RH holds.
\end{corollary}

\begin{proof}
Apply Theorem~\ref{rh-window-input} to $g(u)=G(e^u)$ after the fixed
averaging of Lemma~\ref{rh-averaging}. Formula~\eqref{rh-mellin-formula}
gives meromorphic continuation and simple poles; all real points in
$\Re s>0$ are regular. The growth condition follows from the elementary
bound for $G$. For each $\eta>0$, \eqref{rh-average-target} supplies
arbitrarily large $U=\log Y$ with
$R_U(u)=(\mathcal A_{h_j}G(e^u))_+$ satisfying the required integral
bound. Every detecting pole is preserved. There are consequently no
zeros to the right of $\theta_j$; now let $j\to\infty$.
\end{proof}

It suffices to replace $G$ in \eqref{rh-average-target} by a measurable,
locally integrable upper bound $\mathcal B_j$. For each selected $Y$,
that bound must hold almost everywhere on
$[Y,e^{h_j}Y^{\beta_j}]$, so it covers every argument used by the
local average. Positivity gives
$(\mathcal A_{h_j}G)_+\le(\mathcal A_{h_j}\mathcal B_j)_+$.
With $h_j=0$, a pointwise bound
$G(x)\le M_jx^{\theta_j}$ on arbitrarily late whole intervals gives
an integral at most $M_j(\beta_j-1)\log Y$ and is sufficient.

A signed integral alone does not replace the positive part. For
example, $g(u)=e^{\sigma u}\cos(\tau u)$, with $\theta<\sigma<1$ and
$\tau\ne0$, has a Laplace-transform pole at $\sigma+i\tau$, yet its
signed weighted integral can be negative on arbitrarily late intervals.
It therefore cannot detect the required obstruction by an upper bound
on that signed integral. The positive-part criterion retains it.
